% Généré par scripts/generate_article_tables.py — NE PAS ÉDITER À LA MAIN.
% Source : experiments/exp_S40_article_metrics/summary.json
\begin{table}[htbp]
\centering
\caption{Composantes de la RAM, en octets (\emph{mesuré sur carte}). Le total est la somme des trois premières lignes, le pic de pile étant le maximum des deux phases.}
\label{tab:annex-ram}
\begin{tabular}{lcccc}
\toprule
Composante & \multicolumn{2}{c}{Monitoring} & \multicolumn{2}{c}{Pronostia} \\
\cmidrule(lr){2-3}\cmidrule(lr){4-5}
 & FP32 & INT8 & FP32 & INT8 \\
\midrule
\texttt{.data} & \num{460} & \num{460} & \num{460} & \num{460} \\
\texttt{.bss} & \num{100152} & \num{100152} & \num{105036} & \num{105036} \\
Pic de pile (inférence) & \num{4416} & \num{4328} & \num{4424} & \num{4336} \\
Pic de pile (mise à jour) & \num{4688} & \num{4712} & \num{4696} & \num{4720} \\
Total & \num{105300} & \num{105324} & \num{110192} & \num{110216} \\
\midrule
Rapport & \multicolumn{2}{c}{\num{1.0002}} & \multicolumn{2}{c}{\num{1.0002}} \\
\bottomrule
\end{tabular}
\end{table}
